%% kapa: exam=1 round=35 year=2024 period=1 form=A subject=economics subject_ko=경제학원론 seq=39 number=79
%% source: Q-Net 감정평가사 기출문제 (cst00309) artlSeq=5213229 file=raw/1008/5213229/1교시 문제지 원본.pdf
%% conversion: kordoc 4.24.51 -> tools/gen1.py (auto); review=proofread-sonnet; flags=none
\begin{question}{79}
$t$시점의 실업률은 10\%, 경제활동참가율은 50\%이다. $t$시점과 $t+1$시점 사이에 아래와 같은 변화가 발생할 때, $t+1$시점의 실업률은? (단, 취업자와 비경제활동인구 사이의 이동은 없고, 소수점 둘째자리에서 반올림하여 소수점 첫째자리까지 구한다.)
\begin{framedtext}
○ 실업자 중에서 \par
\hspace*{1em}- 취업에 성공하는 비율(구직률): 20 \% \par
\hspace*{1em}- 구직을 단념하여 비경제활동인구로 편입되는 비율: 10 \% \par
○ 취업자 중에서 실직하여 구직활동을 하는 비율(실직률): 1 \% \par
○ 비경제활동인구 중에서 구직활동을 시작하는 비율: 1 \% \par
\end{framedtext}
\begin{choices}
  \item 8.9 \%
  \item 9.5 \%
  \item 9.9 \%
  \item 10.0 \%
  \item 10.5 \%
\end{choices}
\end{question}
\answer{1}
